\newpage
\section*{Generative AI Tools Usage}
\label{genAIUsage}

An overview of how generative AI tools were used in this thesis is presented in Table \ref{tab:genAIOverview}, and the full details of each usage are given below.

\renewcommand{\arraystretch}{1.3}

\rowcolors{2}{gray!15}{white}
\begin{table}[h!]
    \centering
    \begin{tabularx}{\textwidth}{|
        >{\centering\arraybackslash}m{2.8cm}|
        X|X|X|}
        \hline
        \rowcolor{blue!20}
        AI Tool & Use case & Used in & Remarks \\
        \hline

        \multirow[t]{4}{=}{\centering ChatGPT 5.5}
        & Validation of the bibliography
        & Chapter \ref{cap:relwork} (references)
        & Details below \\
        \cline{2-4}

        & Drafting parts of the written text
        & Sections \ref{sect:bciov}, \ref{sect:limcurrapp},
        \ref{ssect:lsl}, \ref{ssect:pyvizmod}
        & Details below \\
        \cline{2-4}

        & Generation of UML diagrams as PlantUML source code
        & Chapter \ref{cap:sysdes} (figures)
        & Rendered with \url{https://www.planttext.com/} \\
        \cline{2-4}

        & Generation and debugging of application code
        & Application code (\texttt{README.md} and most of the user interface)
        & Details below \\
        \hline
    \end{tabularx}

    \caption{Overview of generative AI usage}
    \label{tab:genAIOverview}
\end{table}

ChatGPT 5.5 was the only generative AI tool used in this thesis. It was used for the four purposes summarized in Table \ref{tab:genAIOverview}. For each of them the purpose of use, the method, the way the output was verified, and the way the output was implemented are described below.

\subsection*{Validation of the bibliography}
\begin{enumerate}
    \item \textbf{Purpose of use:} to check the consistency and plausibility of the bibliographic references collected for the thesis.
    \item \textbf{Method:} the reference list was provided to the model and it was asked to verify that the entries were complete and internally consistent (authors, title, venue, year) and to flag entries that looked incorrect or implausible.
    \item \textbf{Verification of the output:} every reference flagged or suggested by the model was checked manually against the original publication or its official record (publisher page, DOI, or digital library) before being kept. The model was not trusted as a source of citations on its own.
    \item \textbf{Implementation:} only references that were independently confirmed were kept in the bibliography. No citation was added solely on the basis of the model's output.
\end{enumerate}

\subsection*{Drafting parts of the written text}
\begin{enumerate}
    \item \textbf{Purpose of use:} to help draft and rephrase parts of selected sections, namely Section \ref{sect:bciov} (Brain--Computer Interfaces Overview), Section \ref{sect:limcurrapp} (Limitations of Current Approaches), Section \ref{ssect:lsl} (Lab Streaming Layer Integration), and Section \ref{ssect:pyvizmod} (Python Visualization Module).
    \item \textbf{Method:} the model was given the topic, the surrounding context, and the points that each passage had to cover, and was asked to produce draft paragraphs that I then revised.
    \item \textbf{Verification of the output:} all generated text was read critically, checked for technical accuracy against the actual system and the cited literature, and corrected where needed. Any factual claim was confirmed before being kept.
    \item \textbf{Implementation:} the generated drafts were edited and rewritten to match the style and content of the rest of the thesis. The final wording of these sections is my own and reflects the work actually carried out.
\end{enumerate}

\subsection*{Generation of UML diagrams}
\begin{enumerate}
    \item \textbf{Purpose of use:} to obtain the UML diagrams used in the thesis (Chapter \ref{cap:sysdes}).
    \item \textbf{Method:} the model was given a description of the architecture and components and was asked to produce the corresponding diagrams as PlantUML source code.
    \item \textbf{Verification of the output:} the generated PlantUML code was reviewed for correctness against the real system design and adjusted where it did not match the implemented architecture.
    \item \textbf{Implementation:} the reviewed PlantUML code was rendered into the diagram images using \url{https://www.planttext.com/}, and the resulting figures were included in the thesis.
\end{enumerate}

\subsection*{Generation and debugging of application code}
\begin{enumerate}
    \item \textbf{Purpose of use:} to assist with parts of the application code, in particular writing the project \texttt{README.md} and most of the user interface, as well as to help with debugging.
    \item \textbf{Method:} the model was asked to generate UI code and documentation from a description of the desired behaviour, and was given error messages and faulty code snippets to help diagnose problems during development.
    \item \textbf{Verification of the output:} all generated code was reviewed and tested by running the application, and was kept only when it behaved correctly. Suggested fixes were checked before being applied.
    \item \textbf{Implementation:} the verified code was integrated into the application, adapted where necessary to fit the existing code base. The functional logic of the system remains my own work.
\end{enumerate}
